\documentclass[11pt,class=report,crop=false]{standalone}
\usepackage[screen]{../python}


\begin{document}


%====================================================================
\chapitre{Limites}
%====================================================================


\objectifs{La définition de la limite d'une fonction en un point à l'aide de $\epsilon$ et $\alpha$ est la base de toute l'analyse. Dans ce chapitre, on rappelle cette définition, on dégage une méthode pour l'utiliser en pratique, puis on l'applique sur des exercices corrigés.}



%%%%%%%%%%%%%%%%%%%%%%%%%%%%%%%%%%%%%%%%%%%%%%%%%%%%%%%%%%%%%%%%%%%%%
\section{Définition de la limite}

\index{limite}

%--------------------------------------------------------------------
\subsection{Définition}

Soit $f$ une fonction d'ensemble de définition $D_f$, soit $x_0$ un réel et $\ell$ un réel.

\begin{definition}
On dit que $f$ admet pour \defi{limite} $\ell$ en $x_0$, et on note $\displaystyle\lim_{x\to x_0} f(x)=\ell$, si :
\mybox{$(\forall \epsilon>0)\ (\exists \alpha>0)\ (\forall x\in D_f)\quad \big(0<|x-x_0|<\alpha \implies |f(x)-\ell|<\epsilon\big)$.}
\end{definition}

Autrement dit, on peut rendre $f(x)$ aussi proche que l'on veut de $\ell$ (à $\epsilon$ près), à condition de prendre $x$ suffisamment proche de $x_0$ (à $\alpha$ près).

\begin{remarque*}
Le réel $\alpha$ dépend de $\epsilon$ : plus on exige que $f(x)$ soit proche de $\ell$ (c'est-à-dire plus $\epsilon$ est petit), plus il faut en général prendre $x$ proche de $x_0$ (c'est-à-dire $\alpha$ petit).
\end{remarque*}


%--------------------------------------------------------------------
\subsection{Méthode}

Pour montrer $\displaystyle\lim_{x\to x_0} f(x)=\ell$ à l'aide de la définition $(\epsilon,\alpha)$ :
\begin{enumerate}
  \item \textbf{Majoration locale.} On se place sur un voisinage de $x_0$, par exemple $|x-x_0|<r$, et on majore $|f(x)-\ell|$ par une expression du type
  $$|f(x)-\ell| \le M\,|x-x_0|.$$
  Le fait de rester proche de $x_0$ permet de contrôler les facteurs parasites (comme $|2x+5|$ dans l'exercice ci-dessous).

  \item \textbf{Choix de $\alpha$.} Pour avoir $|f(x)-\ell|<\epsilon$, il suffit que $M|x-x_0|<\epsilon$ et que $x$ reste dans le voisinage. On prend donc
  \mybox{$\alpha = \min\left(\dfrac{\epsilon}{M},\ r\right)$.}
\end{enumerate}



%%%%%%%%%%%%%%%%%%%%%%%%%%%%%%%%%%%%%%%%%%%%%%%%%%%%%%%%%%%%%%%%%%%%%
\section{Exercices corrigés}

%--------------------------------------------------------------------
\subsection{Limite d'un polynôme du second degré}

\begin{exercicecours}
Soit la fonction $f : x \mapsto 2x^2+3x+1$.
Montrer en utilisant la définition que :
$$\lim_{x\to 1} f(x) = 6.$$
\end{exercicecours}

\begin{solution}
Montrons que :
$$(\forall \epsilon>0)\ (\exists \alpha>0)\ (\forall x\in D_f)\quad \big(0<|x-1|<\alpha \implies |f(x)-6|<\epsilon\big).$$

\textbf{1) Majoration locale de $|f(x)-6|$.}

Pour $x\in I=\left]\frac12,\frac32\right[$, on a $|x-1|<\frac12$. On factorise :
$$|f(x)-6| = |2x^2+3x-5| = |x-1|\,|2x+5|.$$
Pour $x\in I$ : $\frac12<x<\frac32$, donc $6<2x+5<8$, d'où $|2x+5|<8$ et :
$$|f(x)-6| \le 8\,|x-1|.$$

\textbf{2) Choix de $\alpha$.}

Pour avoir $|f(x)-6|<\epsilon$, il suffit d'avoir $8|x-1|<\epsilon$ et $|x-1|<\frac12$, c'est-à-dire :
$$|x-1|<\frac{\epsilon}{8} \qquad\text{et}\qquad |x-1|<\frac12.$$
Il suffit donc de prendre
$$\alpha=\min\left(\frac{\epsilon}{8},\ \frac12\right).$$

\textbf{Conclusion.} $\displaystyle\lim_{x\to 1} f(x) = 6$.
\end{solution}


\end{document}
